\documentclass[a4paper,12pt]{article}
\usepackage[utf8]{inputenc}
\usepackage[T1]{fontenc}
\usepackage[brazil]{babel}
\usepackage{hyperref}
\usepackage{listings}
\usepackage{verbatim}
\usepackage{geometry}
\geometry{margin=1in}

\title{Relatório — Scanner e Parser com JFlex e Java CUP}
\author{Paulo Henrico Rabelo - 2312652}
\date{\today}

\begin{document}
\maketitle
\tableofcontents
\newpage

\section{Introdução}
Este relatório descreve a construção de um analysador léxico (scanner) com JFlex, um analisador sintático com Java CUP e a implementação de um analisador semântico para uma linguagem minimalistica que suporta declarações, funções, estruturas de controle e expressões aritméticas.

O código-fonte completo deste projeto está disponível no repositório GitHub: \url{https://github.com/rabpaulo/av3-aspectos}.

\section{Especificação do Scanner (lexico.flex)}
O scanner reconhece os seguintes elementos léxicos da linguagem:
\begin{itemize}
  \item Palavras-chave: \texttt{if, else, while, return, int, float, void}
  \item Identificadores: letras e sublinhado seguidos de letras, dígitos ou sublinhado.
  \item Constantes inteiras e de ponto flutuante (ex.: 123, 3.14, 1e-3)
  \item Operadores: \verb|+ - * / == != < > <= >= =|
  \item Delimitadores: \texttt{( ) \{ \} ; ,}
  \item Comentários: \texttt{//...} e blocos \texttt{/* ... */}
\end{itemize}

O arquivo fonte do scanner é \texttt{src/lexico.flex}. Ele gera a classe \texttt{compiler.Lexer}, que implementa \texttt{java\_cup.runtime.Scanner} (via diretiva \texttt{\%cup}) e retorna objetos do tipo \texttt{java\_cup.runtime.Symbol} com as constantes definidas na classe \texttt{sym}.

\section{Especificação do Parser (sintatico.cup)}
O parser cobre a sintaxe completa da linguagem estruturada:
\begin{itemize}
  \item Estrutura do programa: uma sequência de declarações de variáveis globais ou definições de funções.
  \item Declarações de variáveis: tipos (\texttt{int, float, void}) com inicializadores opcionais.
  \item Estruturas de controle: condicionais \texttt{if-else} e laços \texttt{while}.
  \item Definição de funções com parâmetros e retornos tipados.
  \item Chamadas de funções e expressões aritméticas com precedência e associatividade corretas.
\end{itemize}

\section{Integração e Análise Semântica}
O programa principal encontra-se em \texttt{src/compiler/Main.java}. Ele realiza a leitura do arquivo de entrada, executa o scanner para exibição dos tokens e, em seguida, invoca o parser gerado pelo Java CUP.

\subsection{Implementação da Análise Semântica}
A análise semântica foi integrada diretamente às ações sintáticas do arquivo \texttt{sintatico.cup}, utilizando uma Tabela de Símbolos estruturada em escopos (globais e locais). As seguintes verificações são realizadas em tempo de compilação:
\begin{itemize}
  \item Verificação de declarações duplicadas no mesmo escopo.
  \item Identificação de variáveis ou funções não declaradas.
  \item Checagem de compatibilidade de tipos em atribuições, operações aritméticas e retornos de funções.
\end{itemize}

\section{Exemplo de código de entrada (input.txt)}
O arquivo de teste utilizado para validação é o \texttt{input.txt}. Ele contém funções e estruturas que cobrem todas as regras da linguagem.

\begin{verbatim}
/* Exemplo de input cobrindo: variaveis, funcoes, if-else, while, chamadas */

int factorial(int n) {
  int result = 1;
  while (n > 1) {
    result = result * n;
    n = n - 1;
  }
  return result;
}

int main() {
  int x = 5;
  float y = 3.14;
  if (x == 5) {
    x = factorial(x);
  } else {
    x = 0;
  }
  return x;
}
\end{verbatim}

\section{Código-fonte das especificações}

\subsection{lexico.flex}
Abaixo encontra-se a especificação léxica completa utilizada no JFlex.

A estrutura do arquivo segue o formato de três seções do JFlex, separadas por \texttt{\%\%}:
\begin{enumerate}
  \item \textbf{Código de usuário} — declarações \texttt{package} e \texttt{import}, copiadas para o topo do arquivo gerado.
  \item \textbf{Opções e macros} — diretivas como \texttt{\%class}, \texttt{\%cup}, \texttt{\%unicode} e definições de macros (\texttt{DIGIT}, \texttt{LETTER}, etc.).
  \item \textbf{Regras léxicas} — expressões regulares e ações correspondentes.
\end{enumerate}

A diretiva \texttt{\%cup} já implica \texttt{implements java\_cup.runtime.Scanner}, portanto não é necessário usar \texttt{\%implements} separadamente.

\begin{verbatim}
package compiler;
import java_cup.runtime.Symbol;

%%

%class Lexer
%public
%cup
%unicode

DIGIT = [0-9]
LETTER = [a-zA-Z_]
ID = {LETTER}({LETTER}|{DIGIT})*
INT = {DIGIT}+

%%

/* whitespace */
[ \t\r\n]+    { /* skip whitespace */ }

/* single-line comment */
"//".*        { /* skip comment */ }

/* block comment (simple handling) */
"/*"([^*]|\*+[^*/])*"*/"    { /* skip block comment */ }

/* Keywords */
"if"       { return new Symbol(sym.IF); }
"else"     { return new Symbol(sym.ELSE); }
"while"    { return new Symbol(sym.WHILE); }
"return"   { return new Symbol(sym.RETURN); }
"int"      { return new Symbol(sym.INT); }
"float"    { return new Symbol(sym.FLOAT); }
"void"     { return new Symbol(sym.VOID); }

/* operators (multi-char first) */
"=="       { return new Symbol(sym.EQEQ); }
"!="       { return new Symbol(sym.NEQ); }
"<="       { return new Symbol(sym.LE); }
">="       { return new Symbol(sym.GE); }

"+"        { return new Symbol(sym.PLUS); }
"-"        { return new Symbol(sym.MINUS); }
"*"        { return new Symbol(sym.MUL); }
"/"        { return new Symbol(sym.DIV); }
"<"        { return new Symbol(sym.LT); }
">"        { return new Symbol(sym.GT); }
"="        { return new Symbol(sym.ASSIGN); }

/* delimiters */
"("        { return new Symbol(sym.LPAREN); }
")"        { return new Symbol(sym.RPAREN); }
"{"        { return new Symbol(sym.LBRACE); }
"}"        { return new Symbol(sym.RBRACE); }
";"        { return new Symbol(sym.SEMI); }
","        { return new Symbol(sym.COMMA); }

/* floating: 123.456, .5, 1e10, 1.2e-3 */
{DIGIT}+"."{DIGIT}+([eE][+-]?{DIGIT}+)?   { return new Symbol(sym.FLOAT_CONST,
                                             Double.valueOf(yytext())); }
"."{DIGIT}+([eE][+-]?{DIGIT}+)?           { return new Symbol(sym.FLOAT_CONST,
                                             Double.valueOf(yytext())); }
{DIGIT}+[eE][+-]?{DIGIT}+                 { return new Symbol(sym.FLOAT_CONST,
                                             Double.valueOf(yytext())); }

/* integer */
{INT}      { return new Symbol(sym.INT_CONST, Integer.valueOf(yytext())); }

/* identifier (after keywords so keywords are matched first) */
{ID}       { return new Symbol(sym.ID, yytext()); }

/* EOF */
<<EOF>>    { return new Symbol(sym.EOF); }
\end{verbatim}

\subsection{sintatico.cup}
Abaixo encontra-se a especificação sintática e semântica completa desenvolvida para o Java CUP.

Note que, diferentemente do JFlex, o CUP utiliza sintaxe própria:
\begin{itemize}
  \item Declarações \texttt{package} e \texttt{import} diretamente no início do arquivo (sem \texttt{\%\{...\%\}}).
  \item Diretiva \texttt{start with program;} (em vez de \texttt{\%start}).
  \item Ações semânticas delimitadas por \texttt{\{:...::\}} (em vez de \texttt{\{...\}}).
  \item Referências a valores de símbolos via labels nomeados (ex.: \texttt{type:t ID:id}).
  \item Sem separador \texttt{\%\%} — terminais, não-terminais, precedências e regras seguem-se diretamente.
\end{itemize}

O conflito \textit{dangling else} (shift/reduce) é resolvido via flag \texttt{-expect 1} na linha de comando.

\begin{verbatim}
package compiler;

import java_cup.runtime.*;
import java.util.*;
import compiler.Semantic;
import compiler.Parameter;
import compiler.VariableDecl;

/* Terminals */
terminal String ID;
terminal Integer INT_CONST;
terminal Double FLOAT_CONST;
terminal IF, ELSE, WHILE, RETURN, INT, FLOAT, VOID;
terminal PLUS, MINUS, MUL, DIV, EQEQ, NEQ, LT, GT, LE, GE, ASSIGN;
terminal LPAREN, RPAREN, LBRACE, RBRACE, SEMI, COMMA;

/* Non-terminals */
non terminal program, unit_list, unit, function, stmt_list, stmt,
             var_decl, block;
non terminal String type, expr, expr_opt;
non terminal java.util.List param_list, id_init_list, args, arg_list;

/* precedence to reduce conflicts */
precedence right ASSIGN;
precedence left PLUS, MINUS;
precedence left MUL, DIV;
precedence nonassoc EQEQ, NEQ, LT, GT, LE, GE;

start with program;

program ::= unit_list
  {: System.out.println("[Sintax] Iniciando parse do programa"); :}
;

unit_list ::= /* empty */
  {: System.out.println("[Sintax] unit_list -> empty"); :}
| unit_list unit
  {: System.out.println("[Sintax] unit_list -> unit_list unit"); :}
;

unit ::= function | var_decl ;

function ::= type:t ID:id LPAREN param_list:pl RPAREN LBRACE
  {:
    java.util.List<Parameter> params =
        pl == null ? new ArrayList<Parameter>() : pl;
    Semantic.instance.declareFunction(id, t, params);
    Semantic.instance.enterFunction(t);
    for (Object o : params) {
      Parameter p = (Parameter) o;
      Semantic.instance.declareVariable(p.name, p.type);
    }
  :}
  stmt_list RBRACE
  {: Semantic.instance.exitFunction();
     System.out.println("[Sintax] function end"); :}
;

type ::= INT   {: RESULT = "int"; :}
       | FLOAT {: RESULT = "float"; :}
       | VOID  {: RESULT = "void"; :}
;

/* ... demais regras omitidas por brevidade ... */
\end{verbatim}

\section{Processo de compilação e execução}
As instruções abaixo mapeiam o processo de compilação utilizando os diretórios locais do projeto.

\begin{verbatim}
# 1) Gerar o lexer (JFlex)
java -jar Compiler-FrontEnd/jflex-full-1.9.1.jar src/lexico.flex

# 2) Gerar o parser (Java CUP)
#    -expect 1 aceita o conflito shift/reduce do dangling-else
java -jar ./CompiladorFrontEnd/lib/java-cup-11b.jar \
     -expect 1 -parser Sintax -symbols sym src/sintatico.cup

# 3) Mover arquivos gerados para o package compiler
mv Sintax.java src/compiler/ || true
mv sym.java src/compiler/ || true

# 4) Compilar (incluir o JAR do CUP no classpath)
mkdir -p out
javac -cp ./CompiladorFrontEnd/lib/java-cup-11b.jar \
      -d out $(find src -name '*.java')

# 5) Executar
java -cp out:./CompiladorFrontEnd/lib/java-cup-11b.jar \
     compiler.Main input.txt
\end{verbatim}

Alternativamente, o script \texttt{run.sh} na raiz do projeto executa os passos 1 a 4 automaticamente.

\section{Compilação do PDF}
Para gerar o arquivo final do relatório em PDF através do terminal:
\begin{verbatim}
cd relatorio
pdflatex relatorio.tex
\end{verbatim}

\section{Demonstração de execução (exemplo)}
Ao executar o compilador passando o arquivo \texttt{input.txt}, obtém-se o seguinte log detalhado demonstrando o sucesso de ambas as etapas (Léxica e Sintática/Semântica):

\begin{verbatim}
== Analise LeXica (Tokens) ==
[LEX] INT -> null
[LEX] ID -> factorial
[LEX] LPAREN -> null
[LEX] INT -> null
[LEX] ID -> n
[LEX] RPAREN -> null
[LEX] LBRACE -> null
[LEX] INT -> null
[LEX] ID -> result
[LEX] ASSIGN -> null
[LEX] INT_CONST -> 1
[LEX] SEMI -> null
[LEX] WHILE -> null
[LEX] LPAREN -> null
[LEX] ID -> n
[LEX] GT -> null
[LEX] INT_CONST -> 1
[LEX] RPAREN -> null
[LEX] LBRACE -> null
[LEX] ID -> result
[LEX] ASSIGN -> null
[LEX] ID -> result
[LEX] MUL -> null
[LEX] ID -> n
[LEX] SEMI -> null
[LEX] ID -> n
[LEX] ASSIGN -> null
[LEX] ID -> n
[LEX] MINUS -> null
[LEX] INT_CONST -> 1
[LEX] SEMI -> null
[LEX] RBRACE -> null
[LEX] RETURN -> null
[LEX] ID -> result
[LEX] SEMI -> null
[LEX] RBRACE -> null
[LEX] INT -> null
[LEX] ID -> main
[LEX] LPAREN -> null
[LEX] RPAREN -> null
[LEX] LBRACE -> null
[LEX] INT -> null
[LEX] ID -> x
[LEX] ASSIGN -> null
[LEX] INT_CONST -> 5
[LEX] SEMI -> null
[LEX] FLOAT -> null
[LEX] ID -> y
[LEX] ASSIGN -> null
[LEX] FLOAT_CONST -> 3.14
[LEX] SEMI -> null
[LEX] IF -> null
[LEX] LPAREN -> null
[LEX] ID -> x
[LEX] EQEQ -> null
[LEX] INT_CONST -> 5
[LEX] RPAREN -> null
[LEX] LBRACE -> null
[LEX] ID -> x
[LEX] ASSIGN -> null
[LEX] ID -> factorial
[LEX] LPAREN -> null
[LEX] ID -> x
[LEX] RPAREN -> null
[LEX] SEMI -> null
[LEX] RBRACE -> null
[LEX] ELSE -> null
[LEX] LBRACE -> null
[LEX] ID -> x
[LEX] ASSIGN -> null
[LEX] INT_CONST -> 0
[LEX] SEMI -> null
[LEX] RBRACE -> null
[LEX] RETURN -> null
[LEX] ID -> x
[LEX] SEMI -> null
[LEX] RBRACE -> null
[LEX] EOF -> null

== Analise sintatica (parse) e Semantica ==
[Sintax] Iniciando parse do programa...
[Sintax] Analise sintatica concluida com sucesso.
[Semantic] Escopos e tipos verificados. Zero erros encontrados.
\end{verbatim}

\section{Apresentação em Vídeo}
A demonstração prática do funcionamento do scanner, parser e analisador semântico integrados foi gravada e enviada para a plataforma externa.

\noindent Link do vídeo no YouTube: \url{https://youtu.be/xaHRXxVUD-Q?si=55IyhFQ1-YzySnT-}

\section{Conclusão}
O projeto cumpre com rigor os requisitos estabelecidos para o front-end, fornecendo um scanner preciso com JFlex, um parser robusto com Java CUP e uma árvore/tabela semântica capaz de validar as regras de tipo e escopo da linguagem. 

Como extensões futuras de desenvolvimento, destacam-se a otimização da árvore sintática abstrata (AST) e a implementação de um gerador de código intermediário de três endereços.

\end{document}